
\subsection*{6.4. Démonstration de 6.2.2.}
Il suffit de montrer que, quels que soient \(Y\in\operatorname{Ob}(\Sch/S)\), \(y\in Y\) et \(s\in\operatorname{S\ell}(\underline{n})(Y)\), on peut trouver un «~voisinage étale~» \(Y'\to Y\) de \(y\) tel que la restriction de \(s\) sur \(Y'\) soit un produit de commutateurs de \(\operatorname{G\ell}(\underline{n})(Y')\). La question étant locale, on peut supposer \(Y\) affine d'anneau \(A\), et \(\underline{n}\) constant égal à \(n\) sur \(Y\). Soient \(\underline{p}\) l'idéal premier de \(y\) et \((s_{ij})(i,j=1,\ldots,n)\) la matrice de \(s\). Pour tout \(\lambda\in A\) et \(h, k=1,\ldots,n\) ; \(h\ne k\), nous désignerons par \(c_{h,k,\lambda}\) la matrice \((c_{ij})\) \((i,j=1,\ldots,n)\) telle que \(c_{hk}=\lambda\) et \(c_{ij}=\delta_{ij}\) pour \((i,j)\ne(h,k)\) \((\delta_{ij}=\) symbole de Kronecker). Pour toute matrice \(\underline{a}=(a_{ij})(i,j=1,\ldots,n)\) on obtient \(c_{h,k,\lambda}\underline{a}\) en ajoutant à la \(k\)-ème colonne de \(\underline{a}\) \(\lambda\) fois la \(h\)-ème. De même on obtient la matrice \(\underline{a}c_{h,k,\lambda}\) en ajoutant à la \(h\)-ème ligne de \(\underline{a}\) \(\lambda\) fois la \(k\)-ème. En particulier on a
\[
 \det(c_{h,k,\lambda})=1,
\]
et on voit aussitôt que l'on a
\begin{equation}\tag{6.3.1}
 c_{h,k,\lambda}c_{h,k,\mu}=c_{h,k,\lambda+\mu}
\end{equation}
pour tout \(\lambda,\mu\in A\), de sorte qu'on trouve une représentation
\[
 c_{h,k}:G_a\longrightarrow\operatorname{S\ell}(n)
 \qquad(h,k=1,\ldots,n\ ;\ h\ne k).
\]
Sauf mention expresse du contraire, on va supposer dans ce qui suit le \(\lambda\) intervenant dans \(c_{h,k,\lambda}\) \((\lambda\in A\ ;\ h,k=1,\ldots,n\ ;\ h\ne k)\) inversible.
La démonstration se fait alors en deux pas. On va montrer d'abord que \(s\) est un produit de matrices \(c_{h,k,\lambda}\) dans un voisinage zariskien suffisamment petit de \(y\). Puis on va montrer que les matrices \(c_{h,k,\lambda}\) sont des produits de commutateurs de \(\operatorname{G\ell}(n)\) dans un voisinage étale suffisamment petit de \(y\).

